\chapter{Causal response and the Peierls bracket}\label{ch:06}
How can a force applied for a finite time define a tangent vector to the
\emph{unforced} equations? Produce two responses to the same force: one with
no incoming disturbance, the other with no outgoing disturbance. Their
difference solves the homogeneous equations. The endpoint two-form of
Chapter~\ref{ch:05} tells us which difference is the Hamiltonian vector of
the source observable. This construction takes place near a specified
solution, and its domain matters as much as the differential operator.

\section{Fields, sources, and a causal boundary problem}
Fix a smooth solution $\bar\phi$. A perturbation $u$ is a section of the
linearized field bundle $V$. A smooth bulk source $f$ is a section of the
different bundle $V^\vee=V^*\otimes |\Lambda^{\mathrm{top}}T^*M|$; its value
on $u$ is a density, not another field. The pairing and the linearized Euler
operator are
\[
 \langle f,u\rangle=\int_M f(u),\qquad
 J=\left.D{\cal E}\right|_{\bar\phi}:
           \Gamma(V)\longrightarrow\Gamma(V^\vee).
\]
Here $\Gamma(V)$ means smooth sections of $V$, and the oriented Euler form
is read as a density using the spacetime orientation. $J$ is the Hessian
with the sign of the Euler term in
$\delta S$, not an arbitrarily signed wave operator. At this background the
derivative of an observable $F$ is a source when $\delta F(u)=\langle F',u\rangle$.
A choice of metric and measure may identify field components with source
components; it does not exchange their roles in this formula.

For a first source space take smooth dual densities compactly supported in
the spacetime interior. For a nongauge, normally hyperbolic equation on a
globally hyperbolic spacetime without a wall, the causal Green problem for
such tests is well posed \cite{BGP}; reversing the overall sign of the
operator reverses that of its inverses. With a timelike physical wall one
must supply a well-posed initial--boundary problem instead. Use the
homogeneous \emph{linearized} physical boundary condition for both responses
when that condition is its own adjoint. At a cut still held at fixed $j$ the
corresponding trace variation also vanishes. Vanishing past data for a
retarded solution do not mean imposing zero Cauchy data at both ends of a
slab. The advanced problem has the complementary future data.

There are two distinct adjoint assertions here. For interior test sections,
commuting two variations of the action gives formal symmetry $J^*=J$ by
integration by parts. In a slab containing the physical wall the same
calculation, before discarding endpoints, gives
\begin{equation}
 \langle Jv,u\rangle-\langle Ju,v\rangle
 =\omega_{\Sigma_f}(u,v)-\omega_{\Sigma_i}(u,v)
 \label{eq:06-green}
\end{equation}
for $u,v$ obeying the linearized wall condition, with the corner terms in
$\omega$. Indeed the antisymmetrized second bulk variation is
$\langle Ju,v\rangle-\langle Jv,u\rangle$; the complementary terms are the
two endpoint variations. This identity is for trial perturbations, not just
for solutions of $Ju=Jv=0$. Dropping the wall integral for arbitrary fields
would be false. For a scalar, real Robin data $\partial_n u+ku=0$ with
the same fixed $k$ for both variations give
$u\partial_n v-v\partial_n u=0$ on the wall. This is the extra domain check.
An absorbing wall may require a different adjoint boundary condition;
formal symmetry alone does not prove reciprocity for the chosen propagators.

\section{Inverting the endpoint form}
Suppose the adjoint-domain check, existence, uniqueness, and causal support
hold. Write $G^R f$ and $G^A f$ for the retarded and advanced solutions of
$Jv=f$, supported in the causal future and past of $\operatorname{supp}f$,
respectively. On interior source tests $JG^{R/A}=1$. To use an inverse on
the other side, note that $w-G^R Jw$ is the homogeneous solution with the
initial data of $w$. Thus on a finite slab $G^R Jw=w$ requires zero incoming
Cauchy data, and $G^A Jw=w$ requires zero outgoing Cauchy data. Admissible
test fields supported away from both caps meet both conditions; compactness
on a \emph{closed} slab alone does not. Using \eqref{eq:06-green} with a retarded
solution for $f$ and an advanced one for $g$ gives
$\langle g,G^R f\rangle=\langle f,G^A g\rangle$: before the sources the
retarded endpoint vanishes; after them the advanced endpoint vanishes.
Define
\[
 E=G^R-G^A.
\]
Thus $JEf=0$, $EJw=0$ on the common test-field domain just described, and
$E$ is skew with respect to the field--source pairing. For a compact source support
$K$, the support of $Ef$ is in $J^+(K)\cup J^-(K)$, which meets every Cauchy
surface in a compact set in the globally hyperbolic case without a wall.
One-sided inverses can have larger domains: a source turned on after a fixed
slice and left on forever can generate a retarded solution without admitting
a response with zero data after the source is over. One must specify a common
source space before subtracting the two inverses.

For $Ju=0$, take $v=G^R f$ in \eqref{eq:06-green} and put the final slice
after the source. There $G^A f$ vanishes, and at the initial slice $G^R f$
vanishes. Hence $\omega_{\Sigma_f}(u,Ef)=\langle f,u\rangle$. Since $Ef$
is homogeneous, the equality holds on every slice on which the pairing is
finite. Comparing with $\delta H(u)=\omega(u,X_H)$ gives
\begin{equation}
 X_F=E F',\qquad \{F,G\}_{\mathrm P}=\delta F(X_G)
                 =\langle F',E G'\rangle.
 \label{eq:06-peierls}
\end{equation}
With a source added as $S-\varepsilon G$, the linearized equation reads
$Jv=G'$, so its retarded response minus its advanced response is precisely
$X_G$. With $S+\varepsilon G$ each response would carry a minus sign.
Covectors differing by $Jh$ for $h$ in the common test-field domain give
the same tangent pairing and causal vector. Larger classes of sources need their own tests:
the Green operators must act on them, the dual pairings must exist, and
boundary terms must still satisfy \eqref{eq:06-green}. A force applied to an
opened trace pairs with that trace and drives an inhomogeneous boundary
problem; it is not automatically an interior delta source in the fixed-$j$
problem. Likewise singular functional derivatives may require wavefront
conditions before the pairing of distributions in \eqref{eq:06-peierls}
makes sense.

An oscillator determines all of the signs without appealing to a commutator.
For $S=\tfrac m2\int(\dot q^2-\nu^2q^2)dt$, $m,\nu>0$, one has
$J=-m(\partial_t^2+\nu^2)$. With $s=t-t'$ its causal inverses are
\[
 G^R(s)=-\vartheta(s)\frac{\sin(\nu s)}{m\nu},\qquad
 G^A(s)= \vartheta(-s)\frac{\sin(\nu s)}{m\nu}.
\]
Here $\vartheta$ is the step function. Both kernels are continuous and their
first derivatives jump by $-1/m$ at $s=0$, as required by $JG=\delta(s)$.
It follows that $\{q(t),q(t')\}=E(s)=-\sin(\nu s)/(m\nu)$. Differentiating
in the \emph{second} time gives $\{q(t),p(t')\}=m\partial_{t'}E(t-t')
=\cos(\nu(t-t'))$, hence $\{q,p\}=1$ at equal time. Evaluation at a time
can equivalently be obtained from smooth time smearings in this example.

\section{When the Hessian has gauge directions}
An infinitesimal proper gauge transformation $K\varepsilon$ is tangent to the
allowed fields and lies in the presymplectic kernel; a charged cut action
need not be proper. Because gauge transformations take solutions to solutions,
differentiating their action gives $JK=0$ at a solution. Hence $J$ cannot
have an inverse on all fields. The Batalin--Vilkovisky (BV) formalism
\cite{CMR} introduces ghosts for gauge parameters and dual variables
(antifields) for fields and ghosts. A differential $Q$ on these graded
variables encodes the equations, symmetries and their identities in
$Q^2=0$. At its zero $\bar\Phi$ (the classical solution $\bar\phi$ with
ghosts and antifields set to zero), $D=(DQ)_{\bar\Phi}$ squares to zero.
For an irreducible theory, with no further local gauge redundancies among
parameters in $U$, the relevant degrees of the deformation complex are
\[
 \underbrace{\Gamma(U)}_{-1}\xrightarrow{\ K\ }
 \underbrace{\Gamma(V)}_{0}\xrightarrow{\ J\ }
 \underbrace{\Gamma(V^\vee)}_{1}\xrightarrow{\ K^*\ }
 \underbrace{\Gamma(U^\vee)}_{2}.
\]
Here and below the spaces carry compatible support and boundary conditions.
The compositions vanish: the first says that a gauge variation of a solution
solves the linearized equations, the second is the linearized Noether
identity. Boundary tangency is a further requirement. If boundary conditions
are written $b(\Phi)=0$, then $Q$ acts within this problem only if
$b'_\Phi Q(\Phi)=0$ there. At $Q(\bar\Phi)=0$, differentiating in a tangent
direction $w$ gives $b'_{\bar\Phi}Dw=0$. This is why one may restrict the
complex to that boundary domain, rather than testing only its bulk equations.
For Maxwell, $K\varepsilon=d\varepsilon$; fixing the tangential potential on $B$ requires
$d_B(\varepsilon|_B)=0$, which the proper choice $\varepsilon|_B=0$ satisfies.
On a sewn interface matching the two parameter traces similarly preserves
matching of tangential potentials. For a full complex, each arrow must
preserve its specified boundary or transmission domain, not just the field
block in this example.

A \emph{retarded or advanced Green homotopy} $h^{R/A}$ lowers degree by one
and satisfies $Dh^{R/A}+h^{R/A}D=1$ on its corresponding causal support
complex. It is an inverse up to the neighboring gauge and identity arrows.
To construct one, let $Cu=0$ be a linear gauge condition, with $C$ mapping
field perturbations to parameter sections (for Maxwell one can take the
divergence). Choose a symmetric invertible bundle map $B:U\to U^\vee$;
this algebraic weight pairs gauge-condition violations. The gauge-fixed
field operator, the operator giving the change of the gauge condition along
a gauge orbit, and its formal adjoint are respectively
\[
 P=J+C^*BC,\qquad R=CK,\qquad R^*=K^*C^*.
\]
For either causal sign, suppress its superscript and assume $P,R,R^*$
have Green inverses $G_P,G_R,G_{R^*}$ on the corresponding one-sided
domains. The compositions below must preserve support, the field and proper
ghost boundary conditions, and any transmission and homogeneous cap data.
These are hypotheses on the gauge fixing, not consequences of adding a
quadratic term to the action \cite{Khavkine}.

Since $JK=0$, one has $PK=C^*BR$. Therefore $KG_R$ and $G_PC^*B$ solve
the same $P$ problem with the same causal data and are equal by uniqueness.
Likewise $K^*P=R^*BC$ gives $BCG_P=G_{R^*}K^*$. Use the following blocks
on inputs of degrees $0,1,2$, respectively:
\[
 h_0=G_R C,\qquad h_1=G_P,\qquad h_2=C^*G_{R^*}.
\]
For example, on a field perturbation the degree-zero homotopy equation is
$Kh_0+h_1J=KG_R C+G_PJ=G_P(C^*BC+J)=1$. On a source the second
intertwining relation gives $Jh_1+h_2K^*=(J+C^*BC)G_P=1$; the end blocks
are $h_0K=G_RR=1$ and $K^*h_2=R^*G_{R^*}=1$ on their respective domains.

For a concrete check, identify forms with their density duals using the metric on a boundaryless globally hyperbolic spacetime. Denote the formal adjoint of the spacetime exterior derivative by $d^*$, not by the field variation $\delta$. Maxwell gives $K=d$, $J=-d^*d$ on one-forms. Take $C=d^*$ and $B=-1$. With $W_k=dd^*+d^*d$ on $k$-forms, the field block is $P=-W_1$ and the scalar ghost block is $R=W_0$. If $S_k$ is the same-sign causal inverse of $W_k$, the three homotopy blocks are $S_0 d^*$, $-S_1$, and $d S_0$. For instance $d S_0d^*+S_1d^*d=S_1(dd^*+d^*d)=1$: the equality $d S_0=S_1 d$ follows from the differential-operator identity $W_1 d=d W_0$ and causal uniqueness. Adding a wall would oblige us to verify that $d$ really maps the allowed scalar boundary domain into the one-form domain. Commuting the formal differential expressions would not perform that test.

Restoring the retarded and advanced signs and subtracting their identities gives
$D(h^R-h^A)+(h^R-h^A)D=0$. For a closed source $f$ of degree one,
$K^*f=0$, the identity says $Jh^{R/A}f=f$. Thus
$(h^R-h^A)f$ is a homogeneous perturbation. Why should two such gauge choices
give the same bracket? For either causal sign set $w=(h-\widehat h)f$.
Then $Dw=0$. If both homotopies act on the same one-sided support complexes,
with proper parameter domains, their identity on $w$ gives
$w=D\widehat h w=K(\widehat h w)$. This is an allowed proper gauge change.
Pairing it with another closed source $g$ gives zero:
$\langle g,K\varepsilon\rangle=\langle K^*g,\varepsilon\rangle=0$, including
the boundary-domain check. For instance, in Maxwell theory this is the
familiar current conservation test: smearing $d\varepsilon$ against a
compact interior current $j^\mu$ gives
$-\int_M(\nabla_\mu j^\mu)\varepsilon\,|\epsilon|$ after integration by parts.
Therefore the field-block pairing of the causal difference is gauge
independent for the stated class. Skewness uses the adjoint relation of the
two causal field blocks on these sources, as in \eqref{eq:06-green}. A source
with an uncancelled surface term at the cut is not covered by this proof.
If the homotopy identity has right-hand side $1-P_{\mathrm{res}}$ for a
projection onto residual modes, those modes must be retained or treated
separately before making the same inference.

Finally, \eqref{eq:06-peierls} at one background gives a fixed bilinear
bracket for linear observables on the tangent theory. In a nonlinear theory
$J_\phi$, the Green operators, and possibly their domains change with the
solution. Antisymmetry at each background is not the nonlinear Jacobi
identity. One route to that identity requires a smooth (locally reduced)
phase space, Hamiltonian vectors for the admissible observables, and closure
of their algebra: the Jacobi identity then follows from $\delta\omega=0$.
The existence of one causal inverse at $\bar\phi$ establishes none of those
smooth-family assertions.

\begin{exercise}
On coordinates $(x,y,z)$ define a skew bracket by the Leibniz rule and
$\{x,y\}=1$, $\{y,z\}=y$, $\{z,x\}=0$. Does freezing the coefficients at
each point test Jacobi for the original bracket?
\end{exercise}
\noindent\emph{Answer.} Each frozen matrix gives a constant Poisson bracket.
For the varying bracket the cyclic sum on $(x,y,z)$ is
$\{x,\{y,z\}\}+\{y,\{z,x\}\}+\{z,\{x,y\}\}=1$. The missing operation in
the frozen test is the derivative of the bracket itself.
